\documentclass{beamer}

\mode<presentation>
{
  \usetheme{Singapore}      % or try Darmstadt, Madrid, Warsaw, ...
  \usecolortheme{crane} % or try albatross, beaver, crane, ...
  \usefonttheme{serif}  % or try serif, structurebold, ...
  \setbeamertemplate{navigation symbols}{}
  \setbeamertemplate{caption}[numbered]
} 

\usepackage[autostyle=true]{csquotes}
\usepackage{xeCJK}
\usepackage[T1]{fontenc}
\usepackage{xcolor}
\usepackage{setspace}
\setmainfont[AutoFakeSlant=0.2]{Charis SIL}
\setCJKmainfont[AutoFakeSlant=0.2]{Noto Serif CJK HK}
\setCJKmonofont[AutoFakeSlant=0.2]{Noto Serif CJK HK}
\usepackage{polyglossia}
\newfontfamily\thaifont[Script=Thai]{Noto Serif Thai}
\setmainlanguage{english}
\setotherlanguages{thai}
\usepackage{csquotes}
\usepackage{array}
\usepackage{makecell}
\usepackage{changepage}
\usepackage{textcomp}
\usepackage{multirow}

\title[Pearl River Sinitic]{Yuè \& the Sinitic Languages of the Pearl River Region --- Unit 5.1}
\subtitle{13ᵗʰ EACL Summer School, Universität Stuttgart}
\author{Hilário de Sousa 蘇沙}
\institute{EHESS–CRLAO \\ {\tiny membre associé}}
\date{Friday 31 July 2026 \\ 丙午六月十八 \\ {\footnotesize \texttt{hilario@bambooradical.com}}}

\begin{document}

\begin{frame}
  \titlepage
\end{frame}

\section{Hakka (\& Shē)}

\begin{frame}[fragile]{Hakka}

\centering
\includegraphics[width=0.9\linewidth]{lacsechinab8}\\
{\small (客家 Hakka as autonym ⊂ what many linguists consider to be 客家)} 
\end{frame}

\begin{frame}[fragile]{Hakka \& Shē}

\begin{columns}
\begin{column}{0.7\textwidth}
{\footnotesize Coblin (2025: v)'s model}
\includegraphics[width=1\linewidth]{coblinhakkashe}
{\footnotesize
Neo-Hakka 新客 ⊃ people with 客家 identity\\
Shē language 畲話: Sinitic language close to Hakka spoken by people of the 畲 ethnic group\\
}
\end{column}
\begin{column}{0.3\textwidth}
{\footnotesize
Paleo-Hakka 老客 = earlier emigrants, people not identifying as 客家, with partial Neo-Hakka traits (some of these classified are as \enquote{Hakka} in the LAC)\\
\bigskip
Hakka and Shē roughly in complementary distribution (Lo 2015)\\
}
\end{column}
\end{columns}

\vfill \setstretch{0.5} {\tiny Coblin, W. South. 2025. \textit{Common Shē and Common Hakka-Shē: A comparative study} (Language and Linguistics Monograph Series 68). Taipei: Institute of Linguistics, Academia Sinica.\\
Lo, Seo-Gim 羅肇錦. 2015. 畲客關係的名實論證－
客的前世是畲 畲的今生是客. \textit{Journal of Hakka Studies} 客家研究 8(2). 1–26.}\par

\end{frame}

\begin{frame}[fragile]{Hakka \& Shē}

Neo-Hakka 新客 dialects, according to Coblin (2025):\\
1. MC sonorant initial tone B syllables are split into two groups\\
2. \enquote*{to be} 係\\
3. \enquote*{I} *ŋai\\
4. \enquote*{to eat} 食\\
\medskip
Shē 畲：\\
Tone C1 merged into A1\\
Like Neo-Hakka, except \enquote*{to be} 是\\
\medskip
Shē often a little bit influenced by neighbours,\\
e.g. in NE Fujian, \enquote{I} is often ŋwaj (cf. 福州 ŋwaj˧\textsuperscript{B})\\
\medskip
畲 Mand. shē\textsuperscript{A1}, Hakka ciā\textsuperscript{A2}, Canto. se4\textsuperscript{A2}\\ 
\enquote*{slash-and-burn farming}

\end{frame}

\begin{frame}[fragile]{Shē 畲}

\begin{columns}
\begin{column}{0.6\textwidth}
\includegraphics[width=1\linewidth]{sheshe}
\end{column}
\begin{column}{0.4\textwidth}
{\small
「畲族」Shē ethnic group\\
\medskip
「畲話」/山客話: Sinitic language spoken by a few hundred thousand people in the 畲族 ethnic category.\\
\medskip
「畲語」 (shee1238): divergent Hmongic language spoken by <1000 in Guangdong (增城、博羅、惠東、海豐)\\
\smallskip
{\scriptsize (No strong evidence that these two groups are related)}\\
\smallskip
(In Guizhou, another 40000 「東家人」 lumped into the 畲族 category; the speak a type of Western Hmongic)}
\end{column}
\end{columns}

\end{frame}

\begin{frame}[fragile]{Hakka}

\begin{columns}
\begin{column}{0.8\textwidth}
\includegraphics[width=1\linewidth]{watershed2}
\end{column}
\begin{column}{0.25\textwidth}
{\footnotesize
長汀 on 汀江：\\「客家首府」\\ {\tiny highest diversity}
\medskip
梅州 on 梅江：\\「世界客都」\\
\medskip
贛州 on 贛江：\\「客家搖籃」\\
\medskip
東江 area:\\ 
河源「客家古邑」\\
惠州「客家僑都」\\
惠陽\\

\medskip
Huge expansion of Hakka population into e.g. Pearl River Region
}
\end{column}
\end{columns}

\end{frame}

\begin{frame}[fragile]{Hakka}

Original Sinitic inhabitants of Hong Kong:\\
about half 圍頭 Waitau (莞寶粵)\\ about half 客家 Hakka (惠陽腔客家)\\ 
(plus other smaller Yue \& Hakka \& Min varieties)\\
\medskip
Hong Kong Hakka (e.g. Lau 2021)\\ 
{\scriptsize(劉鎮發 Lau Chunfat lau4 \underline{zan3\textsuperscript{C1}} faat3/ Liu Zinfat līu \underline{zǐn\textsuperscript{B}} fǎt)}\\

\vfill {\tiny 
  Lau, Chun Fat 劉鎮發. 2021. 香港客家話研究. Hong Kong: Chung Hwa Book Co. 中華教育.
}

\end{frame}

\section{Initials}

\begin{frame}[fragile]{HK Hakka initials}

b- p- m- f- v-, d- t- l-, z- c- s- j-, g- k- ng- h- (< 劉 Romanisation)\\ (ng(i)- [ȵ], 0- can have [ʔ])\\
\smallskip
(It seems that younger speakers say [ŋi])\\

\bigskip
\begin{tabular}{c|c c c c c}
 & A & \multicolumn{2}{c}{B} & C & D\\ \hline
 & 平 & 被 & 部 & 病 & 白\\
MC & bjæŋ\textsuperscript{A} & bje\textsuperscript{B} & bu\textsuperscript{B} & bjæŋ\textsuperscript{C} & bæk\textsuperscript{D}\\ \hline
香港客 & pʰjaŋ˩\textsuperscript{A2} & pʰi˨˧\textsuperscript{A1} & pʰu˥˧\textsuperscript{C} & pʰjaŋ˥˧\textsuperscript{C} & pʰak˥˧\textsuperscript{D2}\\ \cline{4-6}
Canto. & pʰeŋ˩\textsuperscript{A2} & \multicolumn{1}{c|}{pʰej˩˧\textsuperscript{B2}} & pow˨\textsuperscript{C2} & pɛŋ˨\textsuperscript{C2} & pak˨\textsuperscript{D2}\\ \cline{2-3}
封川粵 & peŋ˨˦˧\textsuperscript{A2} & pi˨˨˧\textsuperscript{B2} & pu˨˩\textsuperscript{C2} & peŋ˨˩\textsuperscript{C2} & pak˧˩\textsuperscript{D2}\\
南寧平 & pən˨˩\textsuperscript{A2} & pəj˩˧\textsuperscript{B2} & pəw˨\textsuperscript{C2} & pən˨\textsuperscript{C2} & pɛk˨\textsuperscript{D2o}\\
漢越 & ɓɪŋ̟˨˩\textsuperscript{A2} & ɓi˧˩ˀ\textsuperscript{B2} & ɓo˧˩ˀ\textsuperscript{B2} & ɓəjŋ̟˧˩ˀ\textsuperscript{B2} & ɓajk̟˧˩ˀ\textsuperscript{D2}\\
\end{tabular}

\vfill {\footnotesize 粵平: Some exceptions e.g. 賊 Canto. caak6\textsuperscript{D2} 南寧平 tʃʰɐk˨\textsuperscript{D2o} \enquote*{thief}\\
客家: Some exceptions e.g. Hakka 渠 gī\textsuperscript{A2} \textsc{3sg} (Canto. 佢  keoi5\textsuperscript{B2})\\
漢越: Except 常船 MC dʑjaŋ\textsuperscript{A} ʑwjen\textsuperscript{A} > SV thường thuyền}\par   
\end{frame}

\section{Finals}

\begin{frame}[fragile]{HK Hakka finals}

\includegraphics[width=0.9\linewidth]{hkhakkafinals}

{\footnotesize 
\begin{enumerate}
    \item The majority of Hakka dialects have lost rounded front vowels
    \item Most HK Hakka speakers have merged -eu into -iu
    \item *-ing *-ik usually > -in -it (e.g. 食 Canto. sik6; Hakka sìt)
\end{enumerate}}
\end{frame}

\begin{frame}[fragile]{HK Hakka finals}

\begin{tabular}{c|c c|c c}
     & MC & 香港客 & Canto. & 南寧平\\ \hline
    沉 & ɖim\textsuperscript{A} & tsʰim˩\textsuperscript{A2} & tsʰɐm˨˩\textsuperscript{A2} & tʃɐm˨˩\textsuperscript{A2}\\
    潛 & dzjem\textsuperscript{A} & tsʰjam˩\textsuperscript{A2} & tsʰim˨˩\textsuperscript{A2} & tʃim˨˩\textsuperscript{A2}\\
\end{tabular}

\bigskip
MC 通攝一等 (-uʷŋ -əʷŋ) vs. 三等 (-juʷŋ -jəʷŋ)\\
\medskip
\begin{tabular}{c|c|c c|c c}
     & MC & 南寧平 & Canto. & 香港客 & 梅縣客\\ \hline
    聾 & luʷŋ\textsuperscript{A} & lʊŋ˨˩\textsuperscript{A2} & lʊŋ˨˩\textsuperscript{A2} & luŋ˨˧\textsuperscript{A1} & luŋ˦\textsuperscript{A1}\\
    龍 & ljəʷŋ\textsuperscript{A} & lʊŋ˨˩\textsuperscript{A2} & lʊŋ˨˩\textsuperscript{A2} & luŋ˩\textsuperscript{A2} & ljuŋ˩\textsuperscript{A2}\\ \hline
    工 & kuʷŋ\textsuperscript{A} & kʊŋ˥˧\textsuperscript{A1} & kʊŋ˥\textsuperscript{A1} & kuŋ˨˧\textsuperscript{A1} & kuŋ˦\textsuperscript{A1}\\
    弓 & kjuʷŋ\textsuperscript{A} & kʊŋ˥˧\textsuperscript{A1} & kʊŋ˥\textsuperscript{A1} & kjuŋ˨˧\textsuperscript{A1} & kjuŋ˦\textsuperscript{A1}\\ \hline
    紅 & ɦuʷŋ\textsuperscript{A} & hʊŋ˨˩\textsuperscript{A2} & hʊŋ˨˩\textsuperscript{A2} & fuŋ˩\textsuperscript{A2} & fuŋ˩\textsuperscript{A2}\\
    胸 & hjəʷŋ\textsuperscript{A} & hʊŋ˥˧\textsuperscript{A1} & hʊŋ˥\textsuperscript{A1} & hjuŋ˨˧\textsuperscript{A1} & hjuŋ˦\textsuperscript{A1}\\ \hline
    屋 & ʔuʷk & ʊk˧\textsuperscript{D2o} & ʊk˥\textsuperscript{DS1} & vuk˧˩\textsuperscript{D1} & vuk˩\textsuperscript{D1}\\
    郁 & ʔjuʷk & ɲʊk˨˧\textsuperscript{D2s} & jʊk˥\textsuperscript{DS1} & juk˧˩\textsuperscript{D1} & juk˩\textsuperscript{D1}\\
  
\end{tabular}

\end{frame}

\section{Tones}

\begin{frame}[fragile]{HK Hakka tones}

\begin{tabular}{c c c c|c c c}
1\textsuperscript{A1} & 2\textsuperscript{A2} & 3\textsuperscript{B} & 4\textsuperscript{C} & (5/3)\textsuperscript{D1} & (6/4)\textsuperscript{D2} & \\ \cline{1-6}
´ & ¯ & ˇ & \`{} & ˇ & \`{} & {\footnotesize (< 劉 Romanisation)}\\
˩˧ & ˩ & ˧˩ & ˥˧ & ˧˩ & ˥˧ & \\
rise & low & mid & high & mid & high & \\
\end{tabular}

\medskip
{\footnotesize 
Hakka: majority of dialects D1 lower than D2 e.g. HK Hk 七 tsʰit˧˩ 十 sip˥˧\\ Pinghua–Yue: all(?) dialects D1 higher than D2 e.g. Canto. 七 tsʰɐt˥ 十 sɐp˨\\
(MC 七 tsʰit\textsuperscript{D} 十 dzjip\textsuperscript{D})}\\

\bigskip
Tone sandhi:\\ 
˩˧ > ˧˥ / \_{}\_{} low or mid\\
客家 hak˧˩ \underline{ka˩˧}, 客家話 hak˧˩ \underline{ka˩˧} va˥˧, 客家人 hak˧˩ \underline{ka˩˧꜔꜒} ŋin˩\\
\smallskip
˥˧ > ˥ / \_{}\_{} low or mid\\
大溪地 tʰaj˥˧ kʰaj˩˧ \underline{tʰi˥˧}, 大溪地人 tʰaj˥˧ kʰaj˩˧ \underline{tʰi˥˧꜒} ŋin˩\\

\end{frame}

\begin{frame}[fragile]{HK Hakka tones}
Coblin (2025) Neo-Hakka 新客: \\
1. MC sonorant initial tone B syllables are split into two groups\\
\bigskip
{\footnotesize
\begin{tabular}{r|c|c|c|c|c|c|c}
MC & \multicolumn{2}{c|}{A} & \multicolumn{2}{c|}{B} & C & \multicolumn{2}{c}{D}\\ \hline
 *voiceless & \multicolumn{2}{c|}{A1 ˩˧\textsuperscript{1}} & \multicolumn{2}{c|}{B ˧˩\textsuperscript{3}} & \multirow{2}{*}{C ˥˧\textsuperscript{4}} & \multicolumn{2}{c}{D1 ˧˩\textsuperscript{3/5}}\\ \cline{3-4} \cline{8-8}
 *vcd sonorant & \textcolor{red}{A1 ˩˧\textsuperscript{1}} & A2 ˩\textsuperscript{2}& \textcolor{red}{A1 ˩˧\textsuperscript{1}} & B ˧˩\textsuperscript{3}& & \textcolor{red}{D1 ˧˩\textsuperscript{3/5}} & D2 ˥˧\textsuperscript{4/6}\\ \cline{2-2} \cline{5-5} \cline{7-7}
 *vcd obstruent & \multicolumn{2}{c|}{A2 ˩\textsuperscript{2}} & \textcolor{red}{A1 ˩˧\textsuperscript{1}} & \multicolumn{1}{c}{C ˥˧\textsuperscript{4}} & & \multicolumn{2}{c}{D2 ˥˧\textsuperscript{4/6}}\\
\end{tabular}}

\bigskip
{\footnotesize The sonorant cases are said to be related to the voicing of sonorants in Proto-Min (Baxter \& Sagart 2014).}
\bigskip

{\footnotesize
\begin{tabular}{c|c c c|c c|c c|c}
 & MC & 南寧平 & Canto. & 香港客 & 長汀客 & PMin & 邵武閩\\ \hline
六 & l-\textsuperscript{D} &  lʊk˨˧\textsuperscript{D2s} & lʊk˨\textsuperscript{D2} & luk˧˩\textsuperscript{D1} & təɰ˨˦\textsuperscript{A2} & **lh- & su˥˧\textsuperscript{D} & 六\\
肉 & ɲ-\textsuperscript{D} & ɲʊk˨˧\textsuperscript{D2s} & juk˨\textsuperscript{D2} & ŋjuk˧˩\textsuperscript{D1} & njəɰ˨˦\textsuperscript{A2} & **nh- & ny˥˧\textsuperscript{D} & 肉\\ \cline{1-1} \cline{5-9}
綠 & l-\textsuperscript{D} & lʊk˨˧\textsuperscript{D2s} & lʊk˨\textsuperscript{D2} &  luk˥˧\textsuperscript{D2} & təɰ˨˩\textsuperscript{C2} & **l- & ly˧˥\textsuperscript{C2} & 綠\\
玉 & ŋ-\textsuperscript{D} & ɲʊk˨˧\textsuperscript{D2s} & juk˨\textsuperscript{D2} &  ŋjuk˥˧\textsuperscript{D2} & njəɰ˨˩\textsuperscript{C2} & **ŋ- & ny˧˥\textsuperscript{C2} & 肉\\

\end{tabular}}

\vfill \setstretch{0.5} {\tiny Baxter, William H. \& Sagart, Laurent. 2014. \textit{Old Chinese: a new reconstruction}. Oxford: Oxford University Press.}\par
\end{frame}


\begin{frame}[fragile]{HK Hakka tones}

{\footnotesize
\begin{tabular}{r|c|c|c|c|c|c|c}
MC & \multicolumn{2}{c|}{A} & \multicolumn{2}{c|}{B} & C & \multicolumn{2}{c}{D}\\ \hline
 *voiceless & \multicolumn{2}{c|}{A1 ˩˧\textsuperscript{1}} & \multicolumn{2}{c|}{B ˧˩\textsuperscript{3}} & \multirow{2}{*}{C ˥˧\textsuperscript{4}} & \multicolumn{2}{c}{D1 ˧˩\textsuperscript{3/5}}\\ \cline{3-4} \cline{8-8}
 *vcd sonorant & \textcolor{red}{A1 ˩˧\textsuperscript{1}} & A2 ˩\textsuperscript{2}& \textcolor{red}{A1 ˩˧\textsuperscript{1}} & B ˧˩\textsuperscript{3}& & \textcolor{red}{D1 ˧˩\textsuperscript{3/5}} & D2 ˥˧\textsuperscript{4/6}\\ \cline{2-2} \cline{5-5} \cline{7-7}
 *vcd obstruent & \multicolumn{2}{c|}{A2 ˩\textsuperscript{2}} & \textcolor{red}{A1 ˩˧\textsuperscript{1}} & \multicolumn{1}{c}{C ˥˧\textsuperscript{4}} & & \multicolumn{2}{c}{D2 ˥˧\textsuperscript{4/6}}\\
\end{tabular}}

\bigskip
{\footnotesize
\begin{tabular}{c|c c c|c c|c c c|c}
 & MC & 南寧平 & Canto. & 香港客 & 長汀客 & PMin & TW & 建甌閩\\ \hline
買 & m-\textsuperscript{B} & maj˩˧\textsuperscript{B2} & maj˩˧\textsuperscript{B2} & maj˩˧\textsuperscript{A1} & me˧\textsuperscript{A1} & **m- & be˥˩\textsuperscript{B} & maj˨˩\textsuperscript{B} & 買\\
懶 & l-\textsuperscript{B} &  lan˩˧\textsuperscript{B2} & lan˩˧\textsuperscript{B2} & lan˩˧\textsuperscript{A1} & laŋ˧\textsuperscript{A1} & & lan˥˩\textsuperscript{B} & laŋ˨˩\textsuperscript{B} & 懶\\ \cline{1-1} \cline{5-10}
耳 & ɲ-\textsuperscript{B} & ɲi˩˧\textsuperscript{B2} & ji˩˧\textsuperscript{B2} & ŋi˧˩\textsuperscript{B} & ni˦˨\textsuperscript{B} & **ŋh- & hĩ˧\textsuperscript{C2} & nejŋ˦˨\textsuperscript{D2} & 耳\\
瓦 & ŋ-\textsuperscript{B} &  wa˩˧\textsuperscript{B2} & ŋa˩˧\textsuperscript{B2} & ŋa˧˩\textsuperscript{B} & ŋwa˦˨\textsuperscript{B} & **ŋh- &  hja˧\textsuperscript{C2} & ua˦˨\textsuperscript{D2} & 瓦\\
網 & m-\textsuperscript{B} & mʊŋ˩˧\textsuperscript{B2} & mɔŋ˩˧\textsuperscript{B2} & mjɔŋ˧˩\textsuperscript{B} & mjɔŋ˦˨\textsuperscript{B} & & baŋ˧\textsuperscript{C2} & mɔŋ˦˨\textsuperscript{D2} & 網\\ 
\end{tabular}}

\end{frame}

\section{Lexicon \& Grammar}

\begin{frame}[fragile]{HK Hakka lexicon \& grammar}

{\footnotesize
(Hakka and Yue share a lot of lexicon. Some differences:)\\
\medskip
Some typical Hakka lexicon, e.g.:\\
圳 \enquote{field furrow} e.g. 深圳 cím zùn\\
埔 \enquote{little plain} e.g. 大埔 tài pù\\
畲/輋 \enquote{slash \& burn (field)} e.g. 禾輋 vō ciā\\ 
\medskip
Many nouns with non-gendered semantics carry sex-specific suffixes\\
虱嫲 sĭt mā \enquote*{louse}, 舌嫲 sèt mā \enquote*{tongue}, 薑嫲 gióng mā \enquote*{fresh ginger}\\
蝠婆 pìt pō \enquote*{bat}, 襖婆 ău pō \enquote*{cotton-padded jacket}\\
蝦公 hā gúng \enquote*{shrimp/prawn}, 貓公 miàu gúng \enquote*{cat}\\
鼻公 pì gúng \enquote*{nose}, 腳趾公 giŏk zĭ gúng \enquote*{big toe}\\
猴哥 hēu gó \enquote*{monkey}, 罌哥 áng gó \enquote*{bottle}\\
\medskip
Cantonese has far fewer of these, e.g.:\\
蝨乸 sat1 naa2 \enquote*{louse}, 蛤乸 gaap3 naa2 \enquote*{frog}\\
腳趾公 goek3 zi2 gung1 \enquote*{big toe}\\
鷯哥 liu4*1 go1 \enquote*{hill myna}, 鼻哥 bei6 go1 \enquote*{nose}\\ 
膝頭哥 sat1 tau4 go1 \enquote*{knee}\\
}
\end{frame}

\begin{frame}[fragile]{HK Hakka lexicon \& grammar}

{\small
\enquote{What}:\\
Hakka: 乜介 mǎk gài {\footnotesize (嘪 mǎi) (gài: 個； OC *kˤa[r]-s)}\\
Canto.: 乜 mat1, 乜嘢 mat1 je5 {\footnotesize (咩 me1, 咩嘢 me1 je5)}\\
\quad (嘢 je5: \enquote*{thing})\\
南寧平: 那門 na˧ mun˩\\

\bigskip
\enquote{To be}:\\ 
Hakka universally uses 係 (港客 hè)\\ 
Yue uses 係 (Canto. hai6), but transitions to using 是 is western Yue (e.g. 玉林粵 ɕi˨˧, 南寧平 ɬi˨)\\
\bigskip
\enquote{To eat}:\\ 
Hakka universally uses 食 (港客 sìt) (at least Neo-Hakka and Shē)\\
Yue: Cantonese uses 食 (sik6), but this is a recent Hakka influence. Amongst Yue dialects 吃/喫 is more common, e.g. 港圍頭 \ 吃 hɛk˧˥\textsuperscript{D1}, 中山石岐 \ 喫 jak˨\textsuperscript{D2}, 南寧平 \ 吃 hɐt˧\textsuperscript{D1}\\
(Canto. also has 喫 jaak3, but now considered very informal in HK.)\\
}
\end{frame}

\begin{frame}[fragile]{HK Hakka lexicon \& grammar}

\includegraphics[width=1\linewidth]{eatdrinktype}
{\footnotesize
Type I e.g. Stan. Mand. 吃飯 chī fàn ≠ 喝茶 hē chá ≠ 抽煙 chōu yān\\
\quad \textit{mainly north of Yangtze}\\
Type II e.g. HK Hakka 食飯 sìt fàn = 食茶 sìt cā = 食煙 sìt yén {\scriptsize (also bòk yén)}\\
\quad \textit{mainly south of Yangtze}\\
Type III e.g. Canto. 食飯 sik6 faan6 = 食煙 sik6 jin1 ≠ 飲茶 jam2 caa4\\
\quad \textit{e.g. Yuè, Huī, Jìn}\\
Type IV eat = drink ≠ smoke\\
\quad {\scriptsize(Their examples: 儋州客\ 食 sìt = sìt ≠ bòk; 南寧江南平\ 吃 hɐt˧ = hɐt˧ ≠ 燒 ʃiu˥)}\\
Type V e.g. 全州湘\ 吃 tʃʒ˦ ≠ 喝酒 xwo˦ ʧjəw˥˦= 喝煙 xwo˦ iɛ̃ŋ˨˩ {\tiny < HM trait}\\
}
\vfill \setstretch{0.5} {\tiny Lu, Wen \& Ding, Hongdi. 2026. The colexification of eating, drinking and smoking in Sinitic languages. \textit{Open Linguistics} 12(1). 20250083. \url{https://doi.org/10.1515/opli-2025-0083}}\par

\end{frame}

\begin{frame}[fragile]{HK Hakka lexicon \& grammar}

\includegraphics[width=1\linewidth]{eatdrinkdiachronic}\\
\bigskip
{\footnotesize (Wu 2000) Cantonese primarily used 喫 during Qīng dynasty (喫 jaak3 = 吃 hek3). 喫 is still more common than 食 amongst Yuè dialects; 食 is usually used amongst Yuè dialects which are (geographically) closer to Hakka and Mǐn dialects.}\\

\vfill \setstretch{0.5} {\tiny Wǔ, Wēi 伍巍. 2000. 粵語動詞"食"、"喫"研究. In Sin, Chow Yiu 單周堯 \& Luke, Kang Kwong 陸鏡光 (eds.), 第七屆國際粵方言研討會論文集 (Fāngyán 方言 Special Volume), 394–396. Beijing: The Commercial Press 商務印書局.}\par
\end{frame}

\begin{frame}[fragile]{HK Hakka lexicon \& grammar}

{\scriptsize

\begin{tabular}{c|c c c|c c c c}
     & \textsc{1sg} & \textsc{2sg} & \textsc{3sg} & \textsc{1pl.ex} & \textsc{1pl.in} & \textsc{2pl} & \textsc{3pl}\\ \hline
    \multirow{2}{*}{港客} & 厓 & 你 & 佢 & 吾兜 & 厓哋 & 惹兜 & 其兜\\
     & ngāi & ngī & gī & ngá déu & ngāi dí & ngiá déu & giá déu\\ \cline{2-8}
    \multirow{2}{*}{\textsc{poss}} & 吾(嘅) & 惹(嘅) & 其(嘅) & 吾兜嘅 & 厓哋嘅 & 惹兜嘅 & 其兜嘅\\
     & ngá(è) & ngiá(è) & giá(è) & ngádéuè & ngāidíè & ngiádéuè & giádéuè\\ \hline\\
    \multirow{2}{*}{Can.} & 我 & 你 & 佢 & \multicolumn{2}{c}{我哋} & 你哋 & 佢哋\\
     & ngo5 & nei5 & keoi5 & \multicolumn{2}{c}{ngo5 dei6} & nei5 dei6 & keoi5 dei6\\ \cline{2-8}
    \multirow{2}{*}{\textsc{poss}} & 我嘅 & etc. & & \multicolumn{2}{c}{我哋嘅} & \textit{etc.} & \\
     & ngo5 ge3 & \textit{etc.} & & \multicolumn{2}{c}{ngo5 dei6 ge3} & \textit{etc.} & \\ \hline
    \multirow{2}{*}{南寧平} & 我 & 你 & 佢 & 我隊 & 伝隊 & 你隊 & 佢隊\\
     & ŋa˩˧ & nəj˩˧ & kəj˩˧ & ŋa˩˧ tɔj˨ & wɐn˨˩ tɔj˨ & nəj˩˧ tɔj˨ & kəj˩˧ tɔj˨\\ \cline{2-8}
    \multirow{2}{*}{\textsc{poss}} & 我个 & \textit{etc.} & & 我隊个 & \textit{etc.} & & \\
     & ŋa˩˧kə˥ & \textit{etc.} & & ŋa˩˧tɔj˨kə˥ & \textit{etc.} & & \\      
\end{tabular}
}
\end{frame}

\begin{frame}[fragile]{HK Hakka lexicon \& grammar}

{\small
(Examples from 劉 2021)\\

\medskip
Comparative: 過 gò \textsc{surpass}\\ 
瓦煲煲飯好食\underline{過}電飯煲\\
ngǎ-báu báu fàn hǎu-sìt \underline{gò} tièn fàn báu\\
casserole cook rice tasty \underline{\textsc{than}} elec. rice pot\\

\medskip
Object marking [\textsc{om}]: 將 zióng \textsc{hold}\\
厓\underline{將}條鎖匙漏嗨在車裏背添\\
ngāi \underline{zióng} tiāu sǒ-sī lèu hói cói cá dí-bòi tiám\\
\textsc{1sg} \underline{\textsc{om}} \textsc{clf} key left \textsc{pfv} at car inside oh!\\

\medskip
Another construction:\\
Object marking [\textsc{om}]: 摎 láu \textsc{help}\\
佢\underline{摎}吾地址寫錯嗨\\
gī \underline{láu} ngá tì-zǐ siǎ cò hói\\
\textsc{3sg} \underline{\textsc{om}} \textsc{1sg:poss} address write wrong \textsc{pfv}\\

}
\end{frame}

\begin{frame}[fragile]{HK Hakka lexicon \& grammar}

{\small
Recipient: 分 bín \textsc{give}\\
{\footnotesize (分 bún e.g. 分配 búnpùi;\\ 
\quad 分 bún > bín \enquote*{give} under Canto. influence? Older Canto. 畀 bi2)}\\

\smallskip
送本書\underline{分}佢\\
sùng bǔn sú \underline{bín} gī\\
gift \textsc{clf} book \underline{\textsc{dat}} \textsc{3sg}\\

\medskip
Permissive causative: 分 bín \textsc{give}\\
你唔好\underline{分}佢去\\
ngī ḿ-mǎu \underline{bín} gī hì\\
\textsc{2sg} \textsc{prohib} \underline{\textsc{give}} \textsc{3sg} go\\

\medskip
Passive: 分 bín \textsc{give}\\
\underline{分}人冤枉\\
\underline{bín} ngīn yén-vǒng\\ 
\underline{\textsc{pass}} person wrongly.accuse\\ 
}
\end{frame}

\section{Hakka 客家 \& Sleiyip 四邑}

\begin{frame}[fragile]{Hakka 客家 \& Sleiyip 四邑}

Hakka 客家 Kèjiā \& Sleiyip 四邑 Sìyì \\
(e.g. Lau 2002, Lin 2021)\\

\bigskip
遷界令／遷海令 The Great Clearance\\ 
(清 edicts 1661, 62, 64, 79)\\
The entire coast (30-50 里 (≈17–29 km) of land) from Shandong southward had to be cleared of civilians and civilian structures, to stop them supplying the 鄭 Dynasty (1628－1683; 明 loyalist) based in Tainan.\\

\vfill \setstretch{0.5} {\tiny Lau, Chun-Fat 劉鎮發. 2002. 粵北土話跟客語和粵語的關係及其歸屬問題 [Northern Guangdong Tuhua’s relationship with Hakka and Yue and the question on its affiliation]. \textit{New Asia Yearly} 新亞論叢 4. 242–249.\\
Lin, Cheng-hui 林正慧. 2021. 由「客人」到「客家」：嘉應州士子以「客」自我定位的過程與影響 From Kheh to Hakka: The process and
influence of the Hakka self-identification in Jiaying Zhou. \textit{Global Hakka Studies} 全球客家研究 16: 79-120.}\par
\end{frame}

\begin{frame}[fragile]{Hakka 客家 \& Sleiyip 四邑}

{\small 展界令/復界令 End of the ban 1683\\
Not all of the original coastal inhabitants returned; large number of people from 嘉應州 moved to e.g. 潮州府、惠州府、廣州府、肇慶府\\
(btw, 嘉應州 was separate from 潮州府 in 1733)}\\
\centering
\includegraphics[width=0.6\linewidth]{qingguangdong}

\end{frame}

\begin{frame}[fragile]{Hakka 客家 \& Sleiyip 四邑}

First opium war 第一次鴉片戰爭 1839–1842 (caused poverty)\\ \smallskip
Red Turban Rebellion 廣東洪兵起義 1854–1856\\
{\footnotesize (清 utilised existing conflicts between「客」(i.e. 嘉應州 migrants) and「土」(本地 Punti i.e. locals), 清 employed 客 to crush the Punti Red Turbans)}\\ \smallskip 
Punti–Hakka Clan Wars 土客械鬥\\
{\footnotesize especially in 四邑 Sìyì / Szeyup (sei3 jap1) / ɬej˧ jip˥ \enquote{four counties} area (1855–1868)}\\
\smallskip
In the end: 清 government expelled most 客 people out of 四邑; some moved back to 嘉應州, some moved to e.g. Guangxi or E Guangdong {\scriptsize(forming huge enclaves)}, some went overseas. {\scriptsize (Due to the fightings, many 四邑 people also went overseas)}\\
\medskip
The exonym 「客家」 was reclaimed by Hakka people in the Pearl River Delta Region, and this usage gradually spread to, e.g., parts of NE Guangdong / W Fujian / S Jiangxi.\\

\end{frame}

\begin{frame}[fragile]{Hakka 客家 \& Sleiyip 四邑}

\centering
\includegraphics[width=0.7\linewidth]{honzijamdin}
{\tiny \url{https://zhuanlan.zhihu.com/p/1906824130126414290}}\\
Classification in 漢字音典 \ {\tiny \url{https://nk2028.shn.hk/yindian/} \url{https://mcpdict.vear.vip/}}\\
\end{frame}

\begin{frame}[fragile]{Hakka 客家 \& Sleiyip 四邑}

\centering
\includegraphics[width=0.9\linewidth]{langatlaschinaguangdong}\\
{\tiny Language Atlas of China 1\textsuperscript{st} ed.}
\end{frame}

\begin{frame}[fragile]{Hakka 客家 \& Sleiyip 四邑}

Traits of 四邑 that are like 粵:\\

\bigskip
\begin{tabular}{c|c c c c c}
 & A & \multicolumn{2}{c}{B} & C & D\\ \hline
 & 平 & 被 & 部 & 病 & 白\\
MC & bjæŋ\textsuperscript{A} & bje\textsuperscript{B} & bu\textsuperscript{B} & bjæŋ\textsuperscript{C} & bæk\textsuperscript{D}\\ \hline
香港客 & pʰjaŋ˩\textsuperscript{A2} & pʰi˨˧\textsuperscript{A1} & pʰu˥˧\textsuperscript{C} & pʰjaŋ˥˧\textsuperscript{C} & pʰak˥˧\textsuperscript{D2}\\ \cline{4-6}
Canto. & pʰeŋ˩\textsuperscript{A2} & \multicolumn{1}{c|}{pʰej˩˧\textsuperscript{B2}} & pow˨\textsuperscript{C2} & pɛŋ˨\textsuperscript{C2} & pak˨\textsuperscript{D2}\\
\textbf{台山} & pʰjaŋ˨\textsuperscript{A2} & \multicolumn{1}{c|}{pʰej˨˩\textsuperscript{B2}} & pu˧˨\textsuperscript{C2} & pjaŋ˧˨\textsuperscript{C2} & pak˧˨\textsuperscript{D2}\\ \cline{2-3}
封川粵 & peŋ˨˦˧\textsuperscript{A2} & pi˨˨˧\textsuperscript{B2} & pu˨˩\textsuperscript{C2} & peŋ˨˩\textsuperscript{C2} & pak˧˩\textsuperscript{D2}\\
南寧平 & pən˨˩\textsuperscript{A2} & pəj˩˧\textsuperscript{B2} & pəw˨\textsuperscript{C2} & pən˨\textsuperscript{C2} & pɛk˨\textsuperscript{D2o}\\
漢越 & ɓɪŋ̟˨˩\textsuperscript{A2} & ɓi˧˩ˀ\textsuperscript{B2} & ɓo˧˩ˀ\textsuperscript{B2} & ɓəjŋ̟˧˩ˀ\textsuperscript{B2} & ɓajk̟˧˩ˀ\textsuperscript{D2}\\
\end{tabular}

\bigskip
Tons of Yue-like vocabulary, e.g.\\
嘢 \enquote*{thing} Canto. je5\textsuperscript{B2}, 中山石岐 ja˨˦\textsuperscript{B}\\
台山 \ 閒嘢 hɛn˨\textsuperscript{A2} njɛ˧\textsuperscript{AC1} / hɛn˨\textsuperscript{A2} njɛ˧˧˥\textsuperscript{AC1*}\\

\end{frame}

\begin{frame}[fragile]{Hakka 客家 \& Sleiyip 四邑}

Traits of 四邑 that are like 客家 / not-like 粵:\\
{\footnotesize (Only talking about core of 四邑 at e.g. 台山, 開平;\\ places closer to Canton e.g. 新會 have a lot more 粵-traits)}\\

\bigskip
Very largely not gone through Pan-Yue vowel changes:\\
\begin{tabular}{c|c c c |c c|c}
     & MC & 香港客 & \textbf{台山} & Canto. & 南寧平 & \\ \hline
    沉 & ɖim\textsuperscript{A} & tsʰim˩\textsuperscript{A2} & tsʰim˨\textsuperscript{A2} & tsʰɐm˨˩\textsuperscript{A2} & tʃɐm˨˩\textsuperscript{A2} & 沉\\
    軍 & kjun\textsuperscript{A} & kjun˨˧\textsuperscript{A1} & kun˧\textsuperscript{AC1} & kʷɐn˥\textsuperscript{A1} & kʷən˥˧\textsuperscript{A1} & 軍\\ \hline
    潛 & dzjem\textsuperscript{A} & tsʰjam˩\textsuperscript{A2} & tʰjam˨\textsuperscript{A2} & tsʰim˨˩\textsuperscript{A2} & tʃim˨˩\textsuperscript{A2} & 潛\\
\end{tabular}

\bigskip
No vowel length difference:\\
\begin{tabular}{c|c c c|c c|c}
    & MC & 香港客 & \textbf{台山} & Canto. & 南寧平 & \\ \hline
    街 & kɛ(j)\textsuperscript{A} & kaj˨˧\textsuperscript{A1} & kaj˨˩\textsuperscript{*} & kaːj˥\textsuperscript{A1} & kaːj˥˧\textsuperscript{A1} & 街\\
    雞 & kej\textsuperscript{A} & kaj˨˧\textsuperscript{A1} & kaj˧\textsuperscript{AC1} & kɐj˥\textsuperscript{A1} & kɐj˥˧\textsuperscript{A1} & 雞\\ \hline
    耕 & kɛŋ\textsuperscript{A} & kaŋ˨˧\textsuperscript{A1} & kaŋ˧\textsuperscript{AC1} & kaːŋ˥\textsuperscript{A1} & kɛːŋ˥˧\textsuperscript{A1} & 耕\\
    更 & kæŋ\textsuperscript{A} & kaŋ˨˧\textsuperscript{A1} & kaŋ˧˧\textsuperscript{AC1} & kɐŋ˥\textsuperscript{A1} & kɐŋ˥˧\textsuperscript{A1} & 更\\
\end{tabular}

\end{frame}

\begin{frame}[fragile]{Hakka 客家 \& Sleiyip 四邑}

Due to there being no vowel-length contrast in 台山:\\ there is a mix of splitting tone D1 by vowel length\\ 
\quad {\footnotesize from Canto. influence (Canto. DS1 ˥ vs. DL1 ˧)}\\
and not splitting tone D1 by vowel length\\ 
\quad {\footnotesize (like Paleo-Hakka tone D1 ˥)}\\ 
{\small i.e. 台山 has more many tone D morphemes in ˥ than Canto.}\\

\bigskip
\begin{tabular}{c|c|c c|c|c}
     & Canto. & 台山 & 惠州 & 香港客 & \\ \hline
    七 & tsʰɐt˥\textsuperscript{DS1} & tʰit˥ \textsuperscript{D1}& tɕʰit˦˥\textsuperscript{D1} & tsʰit˧˩\textsuperscript{D1} & 七\\ \hline
    八 & paːt˧\textsuperscript{DL1} & pat˧\textsuperscript{DL1} {\scriptsize 捌} & & & \\
     & & pat˥\textsuperscript{D1} {\scriptsize 八卦}& pat˦˥\textsuperscript{D1} & pat˧˩\textsuperscript{D1} & 八\\ \hline
    法 & faːt˧\textsuperscript{DL1} & fat˥\textsuperscript{D1} & fat˦˥\textsuperscript{D1} & fap˧˩\textsuperscript{D1} & 法\\ \hline
    駱 & lɔːk˨\textsuperscript{D2} & lɔk˥\textsuperscript{D1} & lɔk˦˥\textsuperscript{D1} & lɔk˥\textsuperscript{D2} & 駱\\
\end{tabular}

\end{frame}

\begin{frame}[fragile]{Hakka 客家 \& Sleiyip 四邑}

\begin{columns}
    \begin{column}{0.5\textwidth}
        台山\\ \smallskip
        \begin{tabular}{c|c|c|c}
             ˧ & ˥ & ˧ & ˥ {\footnotesize /˧}\\ \hline
             ˨ & ˨˩ {\footnotesize /˧/˧˨} & ˧˨ & ˧˨\\
        \end{tabular}
    \end{column}
    \begin{column}{0.5\textwidth}
        惠州\\ \smallskip
        \begin{tabular}{c|c|c|c}
             ˧ & ˧˥ & ˨˩˧ & ˦˥\\ \hline
             ˩ & {\footnotesize ˧˥/˨˩˧/˧˩} & ˧˩ & ˨˩\\
        \end{tabular}
    \end{column}
\end{columns}

\bigskip
{\footnotesize 台山 also has split of MC *sonorant-initial tone B to a degree:}\\
\smallskip
{\small
\begin{tabular}{c|c c|c c c c}
 & MC & Canto. & 香港客 & 台山 & 惠州 & 韶關向陽\textsuperscript{蝨婆}\\ \hline
買 & m-\textsuperscript{B} & maj˩˧\textsuperscript{B2} & maj˩˧\textsuperscript{A1} & maj˧\textsuperscript{AC1} & maj˨˩˧\textsuperscript{C1} & bwə˦\textsuperscript{C1}\\
懶 & l-\textsuperscript{B} &  lan˩˧\textsuperscript{B2} & lan˩˧\textsuperscript{A1} & lan˧\textsuperscript{AC1} & lan˨˩˧\textsuperscript{C1} & lwə˦\textsuperscript{C1}\\ \cline{1-1} \cline{4-6}
耳 & ɲ-\textsuperscript{B} & ji˩˧\textsuperscript{B2} & ŋi˧˩\textsuperscript{B} & ŋej/ŋi˥\textsuperscript{B} & ȵi˧˥\textsuperscript{B} & \multicolumn{1}{|c}{ɡi˥\textsuperscript{C1}}\\ 
瓦 & ŋ-\textsuperscript{B} & ŋa˩˧\textsuperscript{B2} & ŋa˧˩\textsuperscript{B} & 
ŋa˥\textsuperscript{B1} & ŋa˧˥\textsuperscript{B} & \multicolumn{1}{|c}{ɡa˥\textsuperscript{C1}} \\ \cline{7-7}
網 & m-\textsuperscript{B} & mɔŋ˩˧\textsuperscript{B2} & mjɔŋ˧˩\textsuperscript{B} & mɔŋ˥\textsuperscript{B1} & mɔŋ˧˥\textsuperscript{B} & bɔŋ˨˦\textsuperscript{B}\\
\end{tabular}}

\bigskip
(四邑 has merger of tones A1 C1, like 畲)\\
\medskip
台山「我」 ŋᵘɔj˧ \enquote*{I} ŋᵘɔj˨˩ \enquote*{we}\\
惠州「我」 ŋɔj˨˩˧ \enquote*{I} ŋɔj˧˥ \enquote*{we}\\
{\scriptsize 我 with -j not found in Yue except some neighbouring Yue dialects}

\end{frame}

\end{document}
